\documentclass[10pt,a4paper]{amsart}
\usepackage[margin=19mm]{geometry}
\usepackage{amsmath,amssymb,mathtools}
\usepackage[scr=rsfs,cal=euler]{mathalfa}
\usepackage{xfrac}
\usepackage[unicode,hidelinks,bookmarks=false]{hyperref}
\newcommand{\ie}{i.e.\ }
\newcommand{\cf}{cf.\ }
\newcommand{\resp}{respectively\ }
\newcommand{\Subsection}{\subsection}
\providecommand{\nobreakdash}{\nobreak\hbox{-}}
\setlength{\emergencystretch}{2em}
\multlinegap0pt
\usepackage{xcolor}
\usepackage{amssymb}
\usepackage{tikz}
\newtheorem{theorem}{Theorem}
\newcommand{\Rbool}{{\rank}_\mathbb{B}}
\newcommand{\rank}{\mathop{\mathrm{rank}}}
\title{Mathematical and computational perspectives on the Boolean and binary rank and their relation to the real rank}
\author{Michal Parnas}
\date{Source: arXiv:2601.13900v1 (CC BY 4.0)}
\begin{document}
\maketitle

\noindent\emph{Source and licence.} Mathematical and computational perspectives on the Boolean and binary rank and their relation to the real rank. Version arXiv:2601.13900v1 (CC BY 4.0), distributed by arXiv under the Creative Commons Attribution 4.0 International licence, http://creativecommons.org/licenses/by/4.0/, as recorded in the arXiv licence metadata for this exact version and shown on the abstract page https://arxiv.org/abs/2601.13900v1. Full licence notice: \texttt{licenses/arxiv-2601.13900-CC-BY-4.0.txt}.\par\smallskip

\noindent\textit{Editorial scope.} Counted unit: the article's theorem theo-block (source0.tex lines 2131--2135). The article's complete original proof of it is delivered with it (source0.tex lines 2136--2176, one \texttt{proof} environment of 41 lines). No mathematics is written, repaired or re-derived: the statement and the proof are the article's own lines, in the article's own notation, and no retained line is substituted or re-flowed.  Retained uncounted as the source-local context the proof needs: lines 1116--1119; lines 2056--2129.  Nothing was omitted from any retained span, so the package carries no omission record.  No retained line contains a citation, so the wrapper carries no bibliography and no bibliography entry.  No retained line contains a figure environment, an image-inclusion command or drawing code, so no image is delivered and the package declares no asset dependency.  Editorial formatting only, in addition: an AMS \texttt{amsart} preamble that restates the article's own declarations and macros used by the retained lines, namely the environments \texttt{theorem} on separate counters, together with the standard packages those lines need.  The retained environments print in this document as Theorem 1 in order of appearance, whereas the article prints the same items with the numbering of its own full document.  The complete native source of the article as distributed in the arXiv e-print for this exact version is bundled unchanged as \texttt{sources/193054/source0.tex}.
\begin{theorem}
\label{theo:Alon}
Let $M$ be a $0,1$ matrix of size $n \times n$, $n \geq 2$, with at most $d$ zeros in each column. Then $\Rbool(M) \leq c(d) \log_2 n$ where $c(d) =  3e(d+1)/\log_2 e $.
\end{theorem}

\begin{figure}[htb!]
\centering
$$
M =
\left(
\begin{array}{cc}
1 & 1 \\
0 & 1 \\
\end{array}
\right)
\;\;\;\;\;\;\;\;
M^{\otimes 3} =
\left(
\begin{array}{@{}c|c@{}}
\begin{array}{cc|cc}
    1 & 1 & 1 & 1 \\
    0 & 1 & 0 & 1  \\\hline
    0 & 0 & 1 & 1  \\
    0 & 0 & 0 & 1  \\
\end{array}
&
\begin{array}{cc|cc}
     1 & 1 &  1 & 1\\
     0 & 1 &  0 & 1 \\\hline
      0 & 0 &  1 & 1 \\
     0 & 0 &  0 & 1 \\
\end{array}\\\hline
\begin{array}{cc|cc}
    0 & 0 & 0 & 0  \\
    0 & 0 & 0 & 0  \\\hline
    0 & 0 & 0 & 0  \\
    0 & 0 & 0 & 0  \\
\end{array}
&
\begin{array}{cc|cc}
  1 & \textcolor{red}{\bf 1} &  1 & 1 \\
  0 & 1 &  0 & 1 \\\hline
  0 & 0 &  1 & 1 \\
 0 & 0 &  0 & 1 \\
    \end{array}
\end{array}
\right)
$$
\vspace{2em}
\begingroup
   \renewcommand{\arraystretch}{1.6}
$$
M^{\otimes 3} = = \left(\begin{array}{c|c}
  B_{1,1}  &   B_{1,2}\\
\hline
 B_{2,1}  &  {\bf B_{2,2}}
\end{array}\right)
\;\;\;\;\;\;\;\;\;
B_{2,2} = \left(
\begin{array}{c|c}
  {\bf  B'_{1,1}} & B'_{1,2} \\\hline
  B'_{2,1} & B'_{2,2}
  \end{array}\right)
\;\;\;\;\;\;\;\;\;
B'_{1,1} = \left(
\begin{array}{c|c}
   1 & \textcolor{red}{\bf 1} \\\hline
  0 & 1
  \end{array}\right)
$$
\endgroup
\caption{The Kronecker product $M^{\otimes 3}$ of the matrix $M$.
The red one in entry $(5,6)$, that is in row $5$ and column $6$ of $M^{\otimes 3}$,
can be uniquely represented by two vectors $r = (2,1,1)$ and $c = (2,1,2)$, which represent the blocks to which row $5$ and column $6$ belong to, respectively.
For example, row $5$ belongs to block $B_{2,2}$ of $M^{\otimes 3}$, and in this
block to block $B'_{1,1}$ and finally it is in the first row of block $B'_{1,1}$.
Similarly the vector $c = (2,1,2)$ represents the blocks to which column $6$ belongs.}
\label{fig-kronekerblocks}
\end{figure}

\begin{theorem}
\label{theo-block}
Let $M$ be a $0,1$ matrix of size $n \times n$ with at most $d$ zeros in each column. Then for any $k \geq 1$,
$\Rbool(M^{\otimes k}) \leq c(d,k) \cdot \log n$, where $c(d,k)  = 3k\cdot e^k \cdot (d+1)^k \cdot \log_2 e$.
\end{theorem}

\begin{proof}
If $d = 0$ the theorem clearly holds as the Boolean rank of $M^{\otimes k}$ is  $1$ in this case, and if $k = 1$  the theorem follows from Theorem~\ref{theo:Alon}.
Thus, assume that $d \geq 1$ and $k \geq 2$.
Recall that $M^{\otimes k}$ is a matrix of blocks, where each block contains smaller blocks of a certain structure and so on.
Specifically, since $M^{\otimes k} = M \otimes M^{\otimes k-1}$, then at the topmost level $M^{\otimes k}$ is composed of blocks of $M^{\otimes k-1}$ or all zero blocks of the corresponding size.
Inside each non-zero block of the form $M^{\otimes k-1} = M \otimes M^{\otimes k-2}$ there are blocks of the form $M^{\otimes k-2}$ or smaller all-zero blocks, and so on,
until the innermost blocks are either all-zero or equal to $M$, and inside such a block each element can be described as a block of size $1 \times 1$ which is either a one or a zero.

Therefore, any entry $(i,j)$ of $M^{\otimes k}$ can be represented by two vectors $r = (r_1,r_2,...,r_k)$ and $c = (c_1,c_2,...,c_k)$,
which describe the sequence of blocks to which $(i,j)$ belongs, where each block is described by two indices $r_i,c_i$,
such that $r_i$ and $c_i$ are indices of a row and a column of blocks in that level, respectively.
Thus, $(i,j)$ belongs to block $r_1,c_1$ at the topmost level, and then to block  $r_2,c_2$ inside block $r_1,c_1$ and so on,
until the pair $r_k,c_k$ represents a row and a column index inside the innermost block of $M^{\otimes k}$. See Figure~\ref{fig-kronekerblocks} for an illustration.

Let $R$ be a subset of row indices of $M^{\otimes k}$, where each row $i \in R$ is chosen independently and randomly as follows,
according to the sequence of blocks it is contained in. If row $i$ is represented by the row sequence $(r_1,r_2,...,r_k)$,
then  $r_1$ is chosen with probability $1/(d+1)$, and then  $r_2$ is chosen with probability $1/(d+1)$ and so on.
Thus, the row sequence $(r_1,r_2,...,r_k)$ which uniquely represents row $i$, is chosen with probability $1/(d+1)^k$.

Let $C$ be a subset of column indices $j$ of $M^{\otimes k}$, such that for each $i \in R$ it holds that $M^{\otimes k}_{i,j} = 1$.
Let $S  = R \times C$ be the resulting set of pairs $(i,j)$.
The set $S$ is a monochromatic rectangle of ones in $M^{\otimes k}$, since all column indices in $C$ include only ones in the rows in $R$.

What is the probability that a pair $(i,j)$ for which $M^{\otimes k}_{i,j} = 1$ is included in the rectangle defined by $S$?
The probability that row $i$ is selected initially into $R$  is $1/(d+1)^k$.
The pair $(i,j)$ will be in $S$ if for all $i'\in R$ it holds $M^{\otimes k}_{i',j} = 1$, and then $j$ will be included in $C$ (that is, all rows in $R$ have a $1$ in column $j$).
Thus, $(i,j)$ will be in $S$ if we didn't select any row whose row sequence includes an all zero block which includes column $j$,
and since there are at most $d$ zeros in each column of $M$,  this happens with probability at least:
$$
\frac{1}{(d+1)^k}\left( 1 - \frac{1}{(d+1)}\right)^{k\cdot d} \geq \frac{1}{e^k (d+1)^k}
$$
Now if we choose $t =  c(d,k)\cdot \log_2 n  $ rectangles $S_1,...,S_{t}$ independently at random as described above, where $c(d,k)$ is as stated in the theorem,
the probability that $(i,j)$ is not covered by one of these rectangles is at most:
$$
\left(1- \frac{1}{e^k(d+1)^k} \right)^t \leq e^{-t/(e^k(d+1)^k)} = e^{-3k\log_2 n/\log_2 e } = n^{-3k}.
$$
Since there are at most $n^{2k}$ ones in $M^{\otimes k}$, the probability that one of them is not covered by these $t$ rectangles is at most:
$$n^{2k} \cdot n^{-3k} < 1$$
Hence, there exists at least one choice of rectangles $S_1,...,S_t$ for which all ones of $M^{\otimes k}$ are covered,
and thus, $\Rbool(M) \leq t = c(d,k) \cdot \log_2 n $ as required.
\end{proof}

\end{document}
